\section{Modélisation des problèmes}

Chaque problème est décrit à partir des mêmes composantes : état, actions,
fonction de transition, test de but, coût et heuristique.

\begin{table}[H]
\centering
\caption{Résumé de la modélisation des quatre problèmes.}
\label{tab:modelisation}
\begin{tabularx}{\textwidth}{p{2.4cm}Y Y Y}
\toprule
\textbf{Problème} & \textbf{État} & \textbf{Actions} & \textbf{Heuristique} \\
\midrule
Labyrinthe &
Position courante $(r,c)$ &
Déplacements valides vers les cellules voisines &
Distance de Manhattan jusqu'à l'objectif \\
\addlinespace

Navigation &
Ville courante &
Emprunter une route adjacente &
Distance euclidienne à vol d'oiseau \\
\addlinespace

N-Puzzle &
Configuration des tuiles &
Déplacement de la case vide &
Somme des distances de Manhattan des tuiles \\
\addlinespace

Jeu arithmétique &
Valeurs restantes et expressions associées &
Choisir deux termes et appliquer une opération &
Heuristique spécifique ou nulle selon l'instance \\
\bottomrule
\end{tabularx}
\end{table}

\subsection{Labyrinthe}

Dans le problème du labyrinthe, un état est une position $(r,c)$ dans une
grille. Les actions correspondent aux déplacements vers le haut, le bas, la
gauche et la droite, sous réserve que la cellule cible soit accessible.

Lorsque les déplacements sont orthogonaux, l'heuristique de Manhattan est
définie par :

\[
h(n) = |r_n - r_G| + |c_n - c_G|
\]

Dans le cas d'un labyrinthe pondéré, une version admissible peut être obtenue
en multipliant cette distance par le plus faible coût de déplacement possible.

\subsection{Navigation}

Le problème de navigation est représenté par un graphe pondéré. Chaque ville
possède des coordonnées simplifiées $(x,y)$ et les arêtes représentent les
routes disponibles. Le coût d'une action correspond au poids de la route
empruntée.

L'heuristique utilisée est la distance euclidienne à vol d'oiseau :

\[
h(n) =
\sqrt{(x_n-x_G)^2 + (y_n-y_G)^2}
\]

Pour préserver l'admissibilité, les coûts des routes doivent être cohérents
avec les distances géographiques ou être supérieurs à celles-ci.

\subsection{N-Puzzle}

Un état du N-Puzzle est représenté par un tuple immuable contenant la position
de chaque tuile. Les actions possibles déplacent la case vide vers une position
adjacente. Le coût de chaque déplacement est uniforme.

L'heuristique principale est la somme des distances de Manhattan des tuiles
vers leur position cible :

\[
h(s) = \sum_{t \neq 0}
\left(
|r_t-r_t^\star| + |c_t-c_t^\star|
\right)
\]

La case vide, notée $0$, est exclue de cette somme. Le système vérifie aussi
la solvabilité de l'instance avant de lancer une recherche potentiellement
coûteuse.

\subsection{Jeu arithmétique}

Dans le jeu arithmétique, un état contient les valeurs restantes ainsi que les
expressions permettant de les produire. Une action consiste à sélectionner deux
valeurs et à appliquer une opération parmi $+$, $-$, $\times$ et $\div$.

Les fractions exactes sont utilisées afin d'éviter les erreurs d'arrondi liées
aux nombres à virgule flottante. Les divisions par zéro sont interdites. Une
solution est atteinte lorsqu'une valeur restante est égale à la valeur cible.
